\section{Unknown emissions: identification for the model actually fitted}
\label{s4:unknown-emission}
A known measurement matrix is a valuable baseline, but it does not justify
estimating that matrix from the same unvalidated records. A small controlled
model permits a constructive extension. Consider two states and two signals,
an unknown invertible emission matrix $E$ with columns summing to one,
unknown row-stochastic transitions $T_a$, and an unknown positive initial
law $\mu_0$. Require one invertible baseline transition $T_0$. Across
independent trajectories, actions are assigned independently of states,
signals and previous innovations, with positive probabilities. Signal
draws are conditionally independent given their states and invariant across
actions. Conditional on the current state/action, a binary repayment reward
is independent of signal and next-state draws. These restrictions describe
one initial supported POMDP estimation experiment; they are not assumptions
learned from EM convergence.

Condition on $A_0=A_1=0$. Independent assignment makes this the observation
law under two imposed baseline actions. Conditional on $S_1$, the three
signals $O_0,O_1,O_2$ are independent. Let
\[
 \mu_1=T_0'\mu_0,\qquad
 L=E\operatorname{diag}(\mu_0)T_0\operatorname{diag}(\mu_1)^{-1},
 \qquad R=ET_0'.
\]
The columns of $L$ and $R$ are the conditional past/future signal laws
given $S_1$; both are invertible. For each middle signal $j$, define the
observable matrix $C_j(i,k)=P(O_0=i,O_1=j,O_2=k\mid A_0=A_1=0)$.
Then
\begin{align}
 P_{02}:=\sum_jC_j&=L\operatorname{diag}(\mu_1)R',\nonumber\\
 C_jP_{02}^{-1}&=L\operatorname{diag}(E_{j1},E_{j2})L^{-1}.
 \label{s4:controlled-spectral}
\end{align}
Invertibility of a binary emission matrix means the two eigenvalues for
either informative row differ. Its eigenvectors recover $L$ up to
column scaling/permutation. Normalize each column to sum to one; recover
each row of $E$ as the diagonal of $L^{-1}C_jP_{02}^{-1}L$.
The same permutation applies to every row. Now
$\mu_0=E^{-1}p_{O_0}$ and, for each supported action,
\begin{equation}
 T_a=\operatorname{diag}(\mu_0)^{-1}E^{-1}
 P(O_0,O_1\mid A_0=a)E^{-\prime}.
 \label{s4:unknown-transition-recovery}
\end{equation}
For $b_a(s)=P(R_0=1\mid S_0=s,A_0=a)$, the vector
$k_a(o)=P(O_0=o,R_0=1\mid A_0=a)$ gives
$b_a=\operatorname{diag}(\mu_0)^{-1}E^{-1}k_a$.
Thus the unknown-emission, transition and binary-reward model used below
is identified, up to one common state relabeling. No stationary initial
law was required. The state label can be fixed by ordering distress-signal
probabilities; the monetary policy value itself is permutation invariant.

This proof is a project construction using three observed views.
\citet[Sections 3.1--3.2]{spectral2016} develops a broader controlled
spectral route with action-conditioned views, rank assumptions, a
memoryless logging policy and single-trajectory mixing bounds. Its small
singular values enter the finite-sample requirements explicitly. That is
the relevant comparison to an uncontrolled HMM theorem: the action's
dependence on observations changes the views. Augmenting the state by
the action and asserting generic HMM full rank can instead introduce a
rank-deficient transition matrix.

Three limitations follow from the proof itself. An almost singular $E$,
$T_0$ or $P_{02}$ amplifies sampling error; identification is not a
conditioning guarantee. With more hidden states than signal categories,
this two-view inversion fails, even though a different multi-period theorem
may apply. Finally, conditioning on baseline actions chosen using private
information need not leave the experimental law used in the first line.
A logger-specific joint model or additional experimental/proxy information
is then necessary. Under the special private-state-only logger, known
emissions can sometimes recover $\mu(s)q(a\mid s)$ from the joint matrix
and recover $T_a$ by row normalization. Failure of the naive conditional
likelihood is therefore not a universal impossibility theorem for every
private-logger model.

\section{From an identified law to fitted parameters and value uncertainty}
\label{s4:latent-inference}
In the binary experiment write $\vartheta\in\mathbb R^{11}$ for logits
of the initial state probability, two emission probabilities, four
action-specific transition probabilities and four repayment probabilities.
The data for trajectory $i$ are $o_{i,0:H-1},a_{i,0:H-1},r_{i,0:H-1}$.
Conditional on randomized actions, its likelihood is
\[
 L_i(\vartheta)=\sum_{s_{0:H-1}}\mu(s_0)
 \prod_{t=0}^{H-1}E_{o_t,s_t}
     b_{a_t}(s_t)^{r_t}\{1-b_{a_t}(s_t)\}^{1-r_t}
 \prod_{t=0}^{H-2}T_{a_t}(s_t,s_{t+1}).
\]
All final-date reward and emission terms remain present; no final-date
transition is observed in this experiment. The forward recursion sums
this likelihood in $O(HK^2)$ per trajectory and uses a separate scaling
constant at every date. Backward recursion gives $\gamma_{it}(s)$ and
$\xi_{it}(s,s')$ as above. Their expected counts update each Bernoulli
probability. A probability box $[0.002,0.998]$ gives the exact constrained
complete-data maximizer by clipping the expected-count ratio.

Three starts of scaled EM supply initial fits; a direct L-BFGS-B
likelihood maximization using an analytic score provides a numerical
challenger and refinement. For a transition logit $\vartheta_{as}$, the
trajectory score follows by differentiating the hidden-path sum:
\[
 \partial_{\vartheta_{as}}\log L_i
 =\sum_{t=0}^{H-2}\ind(a_{it}=a)
       \{\xi_{it}(s,1)-T_a(s,1)\gamma_{it}(s)\}.
\]
Emission, initial-law and reward logits give analogous posterior success
counts minus fitted probability times posterior trial counts. This is
also a direct derivation of the conditional complete-data score identity.
Finite differencing this analytic score gives the observed information,
including the dependence of the posterior weights on parameters. Treating
the smoothed states as known would omit that missing-state uncertainty.

The implementation uses a centered logit step of $10^{-4}$ for the
information and value gradient. A separate check differentiates the direct
observed likelihood at step $10^{-5}$ and compares with the posterior-count
score; a hidden-path sum independently checks the forward likelihood.
The finite-difference truncation and rounding errors are numerical errors,
not sampling variation. Increasing state dimension or approaching singular
information requires renewed step-size/conditioning analysis. The reported
local curvature diagnostics do not certify a global optimum.

\subsection{The regular sampling argument and its failures}
Take $N$ independent complete trajectories and fixed $H$. Suppose the true
probabilities are in the interior of the declared compact box, rank and
emission separation hold, labels are fixed, and the information matrix
$I=\E[s_i s_i']=-\E[\partial_\vartheta s_i]$ is nonsingular.
For finite alphabets there are finitely many observed trajectory patterns.
Their empirical frequencies obey a uniform law of large numbers; positivity
bounds the pattern log likelihoods and derivatives. The expected log
likelihood is uniquely maximized at the identified law by nonnegativity
of Kullback--Leibler divergence. A global approximate maximizer is
therefore consistent on the label-ordered compact parameter space.
Taylor expansion of its score yields
\begin{equation}
 \sqrt N(\widehat\vartheta-\vartheta_0)
 =I^{-1}\frac1{\sqrt N}\sum_i s_i(\vartheta_0)+o_p(1).
 \label{s4:latent-mle-clt}
\end{equation}
The iid finite-pattern CLT and the information identity then give
covariance $I^{-1}$. This supplies an argument for the actual sampling
law, not a generic claim about a single dependent series or growing $H$.
Multiple EM starts do not prove global maximization. Boundary solutions,
singular information, optimizer caps and disagreement between starts must
remain reported failures or conditional fits.

For a policy $\pi$ fixed before evaluation, let $V_\pi(\vartheta)$ be
the expected discounted recorded financial reward. For finite alphabets
and fixed horizon, summing over hidden paths and observed histories is a
finite smooth function of the environment probabilities, even if $\pi$
uses a threshold. Its history-to-action map is held fixed. Consequently
\[
 \widehat{\operatorname{Var}}(\widehat V_{\pi_1}-\widehat V_{\pi_0})
 =N^{-1}\nabla(V_{\pi_1}-V_{\pi_0})'
       \widehat I^{-1}\nabla(V_{\pi_1}-V_{\pi_0}).
\]
The contrast gradient retains covariance between the policies. Resampling
whole trajectories and refitting all environment parameters is a bootstrap
alternative in this regular regime. It keeps the policy fixed and repeats
initialization, optimization and label handling. A bootstrap that discards
capped fits or recomputes an optimal policy answers a different question.

In the implemented example, $\pi_1$ assists exactly when the binary signal
indicates distress, while $\pi_0$ always uses action zero. Each repayment
pays 100 currency units and assistance costs 12; $H=8$, $\gamma=0.99$.
For a reactive policy, define
\[
 P_\pi(s,s')=\sum_o E_{os}T_{\pi(o)}(s,s'),\qquad
 r_\pi(s)=\sum_o E_{os}\{100b_{\pi(o)}(s)-12\pi(o)\}.
\]
The exact finite-horizon recursion $v_h=r_\pi+\gamma P_\pi v_{h-1}$,
$v_0=0$, gives $V_\pi=\mu'v_H$. This is a fixed observable policy
implemented on a partially observed environment; it does not claim to
be the optimal belief policy. A policy learned from a fitted filter must
retain that trained filter's history-to-action map during all subsequent
environment refits, or use an explicit separate policy-selection design.

\subsection{Numerical tolerance in the same units as policy value}
Within a basin where the negative average log-likelihood Hessian has
smallest eigenvalue at least $\kappa>0$, an average score residual
$\epsilon_s$ bounds parameter optimization error by
$\epsilon_s/\kappa$. A gradient bound $\|\nabla V\|\le G_V$ then
gives value error at most $G_V\epsilon_s/\kappa$. For asymptotic equivalence
in Equation~\eqref{s4:latent-mle-clt}, require
$\epsilon_s=o(N^{-1/2})$ as well as selection of the correct basin.
The observed eigenvalue is a local diagnostic, not proof of a uniform
population curvature bound. Solver success and small likelihood increments
alone do not establish an acceptable monetary error.

Spectral moment inversion is a meaningful alternative: it has a constructive
population solution and avoids iterative latent-path fitting, but inverses
and eigengaps can amplify noisy empirical moments into negative probabilities.
Constrained likelihood uses all recorded patterns and enforces valid
probabilities at a cost in nonconvex optimization. Spectral initialization,
likelihood refinement and the exact small-state reference therefore serve
different purposes. Gaussian filtering, particle likelihood and variational
inference retain the approximation distinctions in the preceding section;
none inherits this finite-alphabet likelihood theorem without a new argument.

\section{A value-error bound on the actual partially observed process}
\label{s4:pomdp-error}
The fully observed bound in Equation~\eqref{r3:simulation-bound} does not
by itself compare two belief processes. A path-law argument does. Align
the same latent labels and observed alphabet in two controlled models.
Let the initial joint state/signal laws differ by at most $\delta_0$ in
$L_1$. Let the joint one-step kernel for reward, next state and next signal,
given current state and action, differ by at most $\delta_K$ uniformly.
Use the \emph{same} policy kernel on every observed history in both models.
Its state in this argument is the complete history together with the latent
state, not just the current observed signal.

Markov kernels contract $L_1$ distance. Insert the first model's one-step
kernel between the two propagated path laws and apply the triangle
inequality. Induction bounds the path-law distance through reward $t$ by
$\delta_0+(t+1)\delta_K$. If $|R_t|\le R$, integration gives
\begin{equation}
 |V_\pi-\widetilde V_\pi|
 \le R\sum_{t=0}^{H-1}\gamma^t
               \{\delta_0+(t+1)\delta_K\}
 \le \frac{R\delta_0}{1-\gamma}
       +\frac{R\delta_K}{(1-\gamma)^2}.
 \label{s4:pomdp-path-bound}
\end{equation}
For factored transition, reward and emission kernels, telescoping the
product measures bounds $\delta_K$ by the sum of their uniform $L_1$
errors. The bound is often conservative, but it explicitly respects partial
observation and policy information. Comparing different optimal policies
or updating the stored policy with each model does not satisfy its premise.

For numerical evaluation of one fixed model/policy, let the backward value
recursion at time $t$ have sup-norm residual $\epsilon_t$. Applying the
same contraction argument backward bounds initial value error by
$\sum_t\gamma^t\epsilon_t$. An infinite-horizon discounted belief
operator gives $\epsilon/(1-\gamma)$ for a uniform Bellman residual.
Both require the operator to be the one induced by the specified observation
and policy process. A grid-spacing comparison alone is not a certified
residual bound. Model error, sampling error, policy selection and numerical
error are distinct terms in a financial decision report.
